\begin{textsample}{NEG Dim 6 – global\_south\_2024\_03 – Score -1.00 – global\_south\_2024\_03/106927521.txt}  \label{ex:f6_neg_388}
More than 80 \% of the employees at Starbucks are \textbf{Malay}, according to Berjaya Group founder, Vincent Tan.
PETALING JAYA : Present and former employees of Starbucks and McDonald ' s in Malaysia are calling for an end to a boycott of the two outlets, saying their livelihoods have been severely impacted by it.
The boycott, meant to express solidarity with Palestinians following sustained military action by Israel in Gaza since Oct 7, is said to have contributed significantly to losses of RM116.32 million suffered by Starbucks owner Berjaya Corp Bhd ( BFood ) in the last quarter of 2023.
ADVERTISEMENT Meanwhile, in a lawsuit brought against the boycott ' s promoters, McDonald ' s Malaysia claimed it suffered losses of RM6 million, including RM1.5 million to lay off employees.
The lawsuit was withdrawn on Friday following a successful court-sanctioned mediation process.
Earlier this month, Berjaya founder Vincent Tan called for the boycotts to end, saying it was only putting the employment of fellow Malaysians at risk.
He said more than 80 \% of the employees at Starbucks are \textbf{Malay}. @ @ @ @ @ @ @ @ @ @ the boycott has affected them personally.
ADVERTISEMENT Varsity student Erica ( not her real name ) lost her part-time job at Starbucks, her sole source of income at the time, due to plummeting sales at her outlet."The sudden loss of income was really difficult for me to deal with.
The other part-timers and I who were let off were left with no job and no money,"she told FMT.
ADVERTISEMENT To make ends meet, Erica is now forced to take all her classes online while working full-time at another job."My salary as a full-timer now is the same as my part-time salary when I was at Starbucks."The boycotts need to stop because many Malaysians like me are affected and have lost our jobs as a result,"she said.
Sarah ( not her real name ), a part-time Starbucks employee, supports the Palestinian cause, but has continued working for the coffeehouse chain.
She said she has been convinced @ @ @ @ @ @ @ @ @ @ links with Israel."However, because of the boycotts, I am no longer able to actively promote Starbucks products to my friends to avoid backlash,"she said.
In a statement issued last month, Starbucks Malaysia denied contributing any part of its profits to the Jewish state ' s government or the Israeli Defence Forces ( IDF ).
It also said its parent, BFood, is a public-listed Malaysian company and that its outlets are operated by Berjaya Starbucks Coffee Company Sdn Bhd, a licensee of Starbucks Coffee International Inc, a US corporation.
In a recent lawsuit, McDonald ' s Malaysia said in incurred losses of RM1.5 million when laying off employees.
Meanwhile, Farha ( not her real name ), who has worked at a McDonald ' s outlet for close to four years now, said it upsets her to see people being let go due to declining sales."In most outlets, including my own, around 95 \% of the workers are Muslims.
So the boycott has @ @ @ @ @ @ @ @ @ @"she said.
Last October, McDonald ' s Malaysia said the company has been a 100 \% Muslim-owned entity since it was acquired in 2017 by Lionhorn Pte Ltd, part of Saudi Arabia ' s Reza Group.
Despite this, Farha said many \textbf{Malays} still believe that McDonald ' s Malaysia endorses Israel.
As a result, they have shunned members of their own community who work at these outlets.
This, she said, has affected the morale of these workers."Many of us employees feel sad about this whole situation."As a result, the employees at her outlet have begun pooling their money to contribute towards NGOs championing the Palestinian cause.
That is something that motivates us to keep going to work every day, she said.

% matched lemmas: malay
\end{textsample}
